\documentclass{subfiles}
\begin{document}
    A problem arises with sources: the alphabet these use are, in general, 
        not suitable for transmission. Is for this reason that,
        as shown in \Cref{fig:shannonSystem}, we make use of transmitters and receivers.
        Although we describe only the encoding process of a source, 
        the deconding is trivially symmetric.

    Let \(S\) be a source. A code is defined as a map from the source alphabet 
        to a new alphabet, called input alphabet. In \Cref{sec:huffmanCodes,sec:arithmeticEncoding,sec:integersEncoding}
        we consider some particular family of codes, 
        while in the remainder of this section we introduce some classes of code,
        as well as some theoretical tools useful for studing codes.

    \subsubsection{Non-singular codes}
    \subfile{../subsub/1.2.1 - Non-singular codes}

    \subsubsection{Uniquely decodable codes}
    \subfile{../subsub/1.2.2 - Uniquely decodable codes and Sardinas-Patterson}

    \subsubsection{Prefix codes and the Kraft-McMillan inequality}
    \subfile{../subsub/1.2.3 - Prefix codes}
\end{document}
